\documentclass[
  paper=b5j,
  fontsize=8pt,
  jafontsize=8pt,
  line_length=42zw
]{jlreq}

\usepackage[
  twocolumn=false,
  footer=false
]{surikumi}

\pagestyle{empty}

\begin{document}

\MMHeading{確認系記号の統一見本}

\MMConfirmationSubsection*{斜線入り正方形}

\noindent
\MMConfirmationStripedMark
\hspace{4em}
\MMConfirmationStripedMark
\hspace{7em}
\MMConfirmationStripedMark

\MMConfirmationSubsection*{黒塗り正方形}

\noindent
\MMConfirmationFilledMark
\hspace{4em}
\MMConfirmationFilledMark
\hspace{7em}
\MMConfirmationFilledMark

\MMConfirmationSubsection*{斜線入りの「矢印＋注」}

\noindent
\MMConfirmationStripedNoteMark
\hspace{3em}
\MMConfirmationStripedNoteMark
\hspace{6em}
\MMConfirmationStripedNoteMark

\MMConfirmationSubsection*{黒塗りの「矢印＋注」}

\noindent
\MMConfirmationFilledNoteMark
\hspace{3em}
\MMConfirmationFilledNoteMark
\hspace{6em}
\MMConfirmationFilledNoteMark

\MMConfirmationSubsection*{本文中の使用例}

\MMConfirmationGeneralComment{斜線入り正方形は，すべての読者に向けたコメントを表す．}

\MMConfirmationAdvancedComment{黒塗り正方形は，意欲的な読者に向けたコメントを表す．}

\MMConfirmationGeneralNote{斜線入りの「矢印＋注」は，すべての読者に向けた注釈を表す．}

\MMConfirmationAdvancedNote{黒塗りの「矢印＋注」は，意欲的な読者に向けた注釈を表す．}

\end{document}
